% ===================================================================
%  اللوحة نمرة ١٤ — الشارع. --- التجّار.
%
%  Traduction, faite en 2026, en arabe egyptien ecrit, sur l'ido de
%  texto/io. Regles de la colonne : voir 10-lawha-01.tex. Une seule
%  numerotation.
%
%  LA RUE DES COMMERCANTS REVELE UN SECOND SUFFIXE DE METIER, A COTE
%  DU -جي DEJA RELEVE : LE SUFFIXE -اتي.
%
%    مصوّراتي (20)  le photographe
%    ساعاتي (37)   l'horloger
%    نضّاراتي (38)  l'opticien
%
%  Il ne se confond pas avec -جي : celui-la se pose sur l'objet
%  qu'on manie — عربجي la voiture, جزمجي la chaussure —, celui-ci sur
%  ce qu'on fabrique ou repare. Le meme tableau donne les deux, et
%  trois de plus en -جي : الصندوقجي (19) l'emballeur, الكوميسيونجي
%  (55) le voyageur de commerce, المطبعجي l'imprimeur. Avec الترزي
%  (62) le tailleur, venu du turc terzi, cela fait sept noms de
%  metier formes de trois facons dans une seule planche.
%
%  ET L'ARMEE QUI DEFILE PORTE TOUTE LA HIERARCHIE OTTOMANE, celle
%  que l'Egypte a gardee jusqu'au milieu du XXe siecle :
%
%    أورطة (41, 47)  le regiment, le bataillon — du turc orta
%    قائمقام         le colonel
%    بكباشي          le lieutenant-colonel
%    صاغ             le commandant
%    يوزباشي         le capitaine — deja au tableau 10
%    طبّالباشي (42)   le tambour-major — avec le -باشي turc
%
%  LE RESTE DE LA RUE :
%
%    فاترينة        la devanture   de l'italien vetrina
%    تنده (6)       le store, la marquise
%    سبيل (78)      la fontaine publique — l'institution du Caire
%    غوايش          les bracelets
%    ألماظ          le diamant     MSA : ماس
%    فكّة            la monnaie
%    حلواني (31)    le patissier
%    دادة (98)      la bonne d'enfants
%    جوخ (67)       le drap
%    فطيرة          la tourte
%
%  LES QUATRE PIEGES DE POSSESSION DU TABLEAU — la rue et ses bains
%  turcs, le bureau du cousin, le jet de la fontaine, l'imperiale de
%  l'omnibus, le tricycle et son coursier — passent tous sans une
%  reprise. Ils avaient coute cinq reecritures a la colonne
%  cantonaise.
% ===================================================================

%%K t14-apar-1 apar
\VUsaut{4.0mm}
\VUtitre{15.0pt}{60}{اللوحة نمرة ١٤}
\VUsaut{3.0mm}

%%K t14-tit-1 sub
\VUpk{11.4pt}{40}{الشارع. --- التجّار.}
\VUsaut{3.0mm}

%%K t14-01-1 p
1. --- صاحبي يوحنا، اللي ساكن في الريف، جه يقضّي تلات أيام عند عيلتي.
\VUgras{ولحدّ ساعتها} ما كانش عنده فرصة يشوف مدينة كبيرة؛ علشان كده إحنا بنقضّي
كل أيامنا في التمشية، وأنا بشرح له اللي هو ما يعرفوش.

%%K t14-02-1 p
2. --- الأول رحنا نزور \VUgras{المرصد} \textsuperscript{(1)}، اللي عجبه أوي. ولما
كنّا راجعين من \VUgras{الشارع} \textsuperscript{(3)} اللي فيه \VUgras{الحمّام
التركي} \textsuperscript{(2)}، قابلنا \VUgras{جنازة} \textsuperscript{(4)}. وقدام
\VUgras{موكب الجنازة} كان ماشي \VUgras{رجال الدين}، قدام \VUgras{عربية الموتى}،
اللي كان متحطّط فيها \VUgras{التابوت}. وأصحاب كتير كانوا ماشيين مع العيلة
\VUgras{الحزينة}.

%%K t14-02-2 p
وبعد ما الموكب عدّى، عدّينا الشارع قدام \VUgras{محل البيرة} \textsuperscript{(5)}
الكبير، اللي فيه \VUgras{زباين} كتير بيشربوا \VUgras{بيرة} على التراس، تحت
\VUgras{تنده} \textsuperscript{(6)} من قزاز.

%%K t14-03-1 p
3. --- ويوحنا اندهش أوي من \VUgras{الدرع} المحطوط بين محل \VUgras{الخمور}
\textsuperscript{(7)} ومحل \VUgras{الآلات الموسيقية} \textsuperscript{(10)}. وسألني عن
معناه؛ فقلت له إنه بيدلّ على \VUgras{القنصلية} \textsuperscript{(8)} الإنجليزية،
ووريته \VUgras{القنصل} \textsuperscript{(9)} اللي على البلكونة.

%%K t14-tit-2 sub
\VUpk{10.2pt}{40}{جولة على المحلات.}
\VUsaut{3.0mm}

%%K t14-04-1 p
4. --- وطبعًا وقفنا قدام فاترينة \VUgras{الآلات الموسيقية}، و\VUgras{الأرغنات
الصغيرة}، و\VUgras{الأرغنات}، و\VUgras{البيانوهات}؛ وبصّينا على أغلفة النوتة

%%K t14-04-1 p suite
المصوّرة وعلى \VUgras{نوتة} البيانو، واتفرّجنا على \VUgras{القيثارة} الخشب المدهّبة
اللي بتتستعمل يافطة.

%%K t14-05-1 p
5. --- وسحب التراب اللي عملها \VUgras{الكنّاسين} \textsuperscript{(11)}
و\VUgras{عربية الكنس} \textsuperscript{(12)} اضطرّتنا نعدّي بسرعة قدام \VUgras{العفش}
الحلو بتاع \VUgras{محل العفش} \textsuperscript{(13)}، وقدام فاترينة
\VUgras{الفوانيس} و\VUgras{اللمبات} بتاعة \VUgras{السمكري} \textsuperscript{(14)}.

%%K t14-06-1 p
6. --- وكمان في المكان ده اضطرّينا نسيب الرصيف، علشان كان مسدود بطرود، اللي
كانوا بينزّلوها من \VUgras{عربية نقل} \textsuperscript{(16)} علشان يحطّوها في
\VUgras{محل} \textsuperscript{(15)} كان لسه اتأجّر.

%%K t14-06-2 p
وده منعنا نشوف العربيات الحلوة بتاعة \VUgras{صانع العربيات} \textsuperscript{(17)}،
بس ما منعناش نقف نتفرّج على \VUgras{الصندوقجي} \textsuperscript{(19)} وهو بيعمل صندوق
لجاره، \VUgras{بتاع العجل والعربيات} \textsuperscript{(18)}. وبعد كده، علشان الوقت كان
اتأخّر شوية، رجعنا البيت.

%%K t14-tit-3 sub
\VUpk{10.2pt}{40}{لفّة في المدينة.}
\VUsaut{3.0mm}

%%K t14-07-1 p
7. --- وتاني يوم، أمي اقترحت على يوحنا إنه يتصوّر، وبعدين كلّفتني إني أروح معاه
عند \VUgras{المصوّراتي} \textsuperscript{(20)}. وبعد الغدا، دخلنا \VUgras{أتيليه}
الفنّان، اللي اضطرّينا نستنّى فيه مدة طويلة شوية. وعلشان نشغّل نفسنا، بصّينا على
\VUgras{الصور} \textsuperscript{(22)} المعلّقة على الحيطة.

%%K t14-07-2 p
ولما جه دورنا، الشغل ما خدش وقت طويل، وأول ما صاحبي خلص وقوف قدام
\VUgras{الكاميرا} \textsuperscript{(21)}، نزلنا.

%%K t14-08-1 p
8. --- وفي الدور الأول، شفنا من باب \VUgras{الحلّاق} \textsuperscript{(23)}، اللي كان
مفتوح، صبيّه وهو بيقصّ شعر زبون.

%%K t14-08-2 p
ولما وصلنا الشارع، عدّينا قدام \VUgras{محل الدخّان} \textsuperscript{(24)}. ويوحنا
اتسلّى أوي بـ\VUgras{الفانوس} \textsuperscript{(25)} الأحمر وبـ\VUgras{الغليون الكبير}
اللي بيتستعمل \VUgras{يافطة} \textsuperscript{(26)}، لدرجة إنه اصطدم باتنين أساتذة
كبار كانوا بيتكلّموا وهما \VUgras{بيشمّوا سعوط} \textsuperscript{(27)}.

%%K t14-tit-4 sub
\VUpk{10.2pt}{40}{الحلواني.}
\VUsaut{3.0mm}

%%K t14-09-1 p
9. --- و\VUgras{البنكنوت} و\VUgras{الفكّة} من كل البلاد المعروضين في
\VUgras{فاترينة} \VUgras{الصرّاف} \textsuperscript{(29)} شدّوا انتباهنا شوية، بس علشان
كانت ساعة الأكل الخفيف، دخلنا عند \VUgras{الحلواني} \textsuperscript{(31)} من غير ما
نقف أكتر.

%%K t14-10-1 p
10. --- ولما شفنا كل \VUgras{الحاجات الحلوة} اللي في المحل، \VUgras{الكيك، وكعك
العسل، والبسكوت}، وكل أنواع \VUgras{الملبّس}، ما قدرناش نختار بسهولة. وفي الآخر
يوحنا اختار \VUgras{الجاتوه}، وأنا خدت بس \VUgras{تارت كرز} صغيرة، علشان الحلويات
\VUgras{بتسوّس سناني}، وأنا مش عايز خالص أروح كتير عند \VUgras{دكتور السنان}
\textsuperscript{(30)}.

%%K t14-tit-5 sub
\VUpk{10.2pt}{40}{الكتابة على الآلة.}
\VUsaut{3.0mm}

%%K t14-11-1 p
11. --- وعلشان أوري صاحبي \VUgras{الآلة الكاتبة} اللي كان سمع عنها بس ما شافهاش،
دخلنا \VUgras{مكتب} \textsuperscript{(32)} \VUgras{ابن عمّي} \textsuperscript{(34)}.
ولقيناه بيملي جواباته على \VUgras{سكرتيره} \textsuperscript{(35)} اللي بيعرف
\VUgras{الاختزال}. وأول ما الجواب يخلص، \VUgras{الكاتب على الآلة}
\textsuperscript{(33)} كان بينقله على آلته. وعلشان ما نعطّلش شغل قريبي، ما قعدناش
عنده غير كام دقيقة.

%%K t14-12-1 p
12. --- وتحت مكتب ابن عمّي، ساكن \VUgras{ساعاتي وجواهرجي} \textsuperscript{(37)}
كبير، اللي كمان

%%K t14-12-1 p suite
بيشتغل في الدهب. ومحلّه الرائع فيه \VUgras{ساعات} كتير، و\VUgras{ساعات حيطة}،
و\VUgras{ساعات جيب}، و\VUgras{سلاسل}، و\VUgras{مجوهرات} غالية، زي \VUgras{عقود
اللؤلؤ}، و\VUgras{الغوايش} الدهب، و\VUgras{الحلق}، و\VUgras{الخواتم} المرصّعة
بـ\VUgras{الأحجار الكريمة} و\VUgras{الألماظ}. وواحدة من الفاترينات مخصّصة
لـ\VUgras{المصوغات} وفيها مجموعة كبيرة من \VUgras{أطقم} الفضّة.

%%K t14-13-1 p
13. --- ولما لفّينا ناصية الشارع، لاحظت إنهم غيّروا \VUgras{اللافتة}
\textsuperscript{(36)} المحطوطة جنب \VUgras{نمرة} البيت. وكنت هقول ملاحظتي بصوت عالي،
لما سمعنا الأصوات المفرحة بتاعة \VUgras{الموسيقى} العسكرية. وجرينا نستقبل الأورطة،
من غير ما نفكّر إن \VUgras{الأورطة} هتعدّي قدام محلات \VUgras{البرانيط}
\textsuperscript{(39)} و\VUgras{الجوانتي} \textsuperscript{(40)}، وإننا كنّا هنشوف
عرضها كويس زي كده لو فضلنا قدام فاترينة \VUgras{النضّاراتي} \textsuperscript{(38)}،
المليانة

%%K t14-13-1 p suite
\VUgras{نضّارات أنف}، و\VUgras{مونوكلات}، و\VUgras{نضّارات مسرح}.

%%K t14-tit-6 sub
\VUpk{10.2pt}{40}{العرض العسكري.}
\VUsaut{3.0mm}

%%K t14-14-1 p
14. --- قدام \VUgras{الأورطة} \textsuperscript{(41)} كان ماشي \VUgras{الطبّالباشي}
\textsuperscript{(42)}، وراه \VUgras{الطبّالة} \textsuperscript{(43)}،
و\VUgras{البوّاقين} \textsuperscript{(44)}، وفرقة \VUgras{الموسيقى} كلها.
و\VUgras{الجنرال} \textsuperscript{(45)} اللي كان في الميدان مع ياوره، كان واقف جنب
\VUgras{نفّاث المية} \textsuperscript{(49)} بتاع \VUgras{النافورة} \textsuperscript{(48)}
علشان يشوف \VUgras{العرض}.

%%K t14-14-2 p
\VUgras{وضبّاط الأركان} \textsuperscript{(46)}، و\VUgras{القائمقام}،
و\VUgras{البكباشي} سلّموا عليه بالسيف. \VUgras{والصاغات} كانوا قدام
\VUgras{أورطاتهم} \textsuperscript{(47)}، وكل \VUgras{يوزباشي} كان بيقود
\VUgras{سريّته}، ولمّا \VUgras{الملازمين}، و\VUgras{الملازمين التانيين}،
و\VUgras{ضبّاط الصف} كانوا في أماكنهم المعتادة جنب رجّالتهم.

%%K t14-15-1 p
15. --- وناس كتير مارّة اتجمّعوا على \VUgras{التقاطع} \textsuperscript{(50)}؛
والتجّار، \VUgras{وبتاع الشماسي} \textsuperscript{(51)}، و\VUgras{السمكري}
\textsuperscript{(53)}، و\VUgras{بتاع السلاح} \textsuperscript{(54)} خرجوا يشوفوا
\VUgras{العساكر}. وكل العربيات، \VUgras{العربيات الأجرة} \textsuperscript{(52)}،
و\VUgras{عربيات النقل} \textsuperscript{(57)}، وكمان \VUgras{عربية الرشّ}
\textsuperscript{(56)}، ركنت على طول الرصيف.

%%K t14-tit-7 sub
\VUpk{10.2pt}{40}{مناظر من الشارع.}
\VUsaut{3.0mm}

%%K t14-15-2 p
و\VUgras{كوميسيونجي} \textsuperscript{(55)} شايل بإيده \VUgras{علب العيّنات} بتاعته،
عدّى الشارع بسرعة، والناس القاعدين على \VUgras{الطابق العلوي} \textsuperscript{(59)}
بتاع \VUgras{الأومنيبوس} \textsuperscript{(58)} لفّوا يتفرّجوا على المنظر.
و\VUgras{مغنّي متجوّل} \textsuperscript{(60)} مسكين بـ\VUgras{أرغنه اليدوي}
\textsuperscript{(61)} اضطرّ يقطع \VUgras{أغنيته}، علشان دوشة الطبول غطّت على صوته.

%%K t14-16-1 p
16. --- ولما الأورطة وصلت قدام \VUgras{محل الأقمشة} \textsuperscript{(64)}،
\VUgras{الخيّاطات} \textsuperscript{(63)} و\VUgras{الترزية} \textsuperscript{(62)} سابوا
\VUgras{مكن الخياطة} وشغلهم علشان يتفرّجوا من الشباك. و\VUgras{صاحب المحل}
\textsuperscript{(65)}، اللي كان لسه حاطط \VUgras{بدل} \textsuperscript{(66)} جديدة في
\VUgras{الفاترينة}، اضطرّ يدخّل \VUgras{الجوخ} \textsuperscript{(67)}
و\VUgras{الأقمشة} اللي كانت مفرودة على الرصيف، جنب \VUgras{البالوعة}
\textsuperscript{(68)}، خوفًا من إن الزحمة تقلبها؛ وحتى \VUgras{البيّاع المتجوّل}
\textsuperscript{(69)} العجوز، اللي كانت واحدة بوّابة بتفاصله على شرّابات، اضطرّ يدخل
في طرقة علشان يخلّص بيعه.

%%K t14-16-2 p
ومشينا مع الأورطة لحدّ القشلاق، \VUgras{و}إحنا مبسوطين أوي من يومنا، رجعنا في
ساعة العشا.

%%K t14-tit-8 sub
\VUpk{10.2pt}{40}{حرف ومهن مختلفة.}
\VUsaut{3.0mm}

%%K t14-17-1 p
17. --- وتاني يوم، عمّي اقترح علينا نزور مطبعة الجرنال اللي هناك. ووافقنا بحماس.

%%K t14-17-2 p
و\VUgras{المطبعجي} كمان \VUgras{صاحب مكتبة} \textsuperscript{(70)} \VUgras{(يعني بيبيع
كتب)}، و\VUgras{ناشر}، و\VUgras{ورّاق}، و\VUgras{مجلّد}. وهو عامل في الدور الأول
\VUgras{التحرير} \textsuperscript{(71)} بتاع الجرنال، وكمان \VUgras{قسم الحفر} اللي
بيشتغل فيه \VUgras{الحفّارين على الحجر} \textsuperscript{(72)} و\VUgras{الحفّارين على
النحاس} \textsuperscript{(73)}، وأخيرًا \VUgras{قسم الجمع}، اللي فيه \VUgras{عمال
الجمع} \textsuperscript{(74)}. و\VUgras{المطبعة} \textsuperscript{(75)} نفسها،
و\VUgras{مكن الطبع}، والماكينات المختلفة، في الدور الأرضي.

%%K t14-tit-9 sub
\VUpk{10.2pt}{40}{يوحنا بيسأل أسئلة كتير أوي.}
\VUsaut{3.0mm}

%%K t14-18-1 p
18. --- ولما خرجنا من المطبعة، يوحنا وقف مندهش أوي قدام علبة قزاز صغيرة محطوطة
على عمود، قصاد \VUgras{محل الخردوات} \textsuperscript{(77)}، وسأل بتتستعمل في إيه.
وعمّي شرح له إنها \VUgras{جرس الحريقة} \textsuperscript{(76)}. وبالمناسبة كل حاجة في
المدينة كانت بتدهش صاحبي، من \VUgras{السبيل} \textsuperscript{(78)} البسيط
و\VUgras{التريسيكل} \textsuperscript{(80)} الخفيف اللي بيسوقه \VUgras{صبي}
\textsuperscript{(79)} بلبس مميّز، لحدّ \VUgras{عربية البوستة} \textsuperscript{(81)}.

%%K t14-19-1 p
19. --- ولما عمّي سابنا شوية علشان يتكلّم مع \VUgras{محصّل الضرايب}
\textsuperscript{(83)}، اللي كان واقف بيتكلّم مع \VUgras{محامي} \textsuperscript{(82)}
و\VUgras{موثّق} \textsuperscript{(84)}، فضلنا نسلّي نفسنا بمتابعة الحركة في الشارع
بعينينا. وناس من كل نوع عدّت قدامنا: \VUgras{صبي حلواني} \textsuperscript{(85)} شايل
في مقطف \VUgras{فطيرة} رائعة، جه ورا \VUgras{بتاع الهدوم} \textsuperscript{(86)}؛
و\VUgras{بتاع النور} \textsuperscript{(87)} كان بيتكلّم مع \VUgras{بتاع الغاز}
\textsuperscript{(88)}؛ و\VUgras{راهبة} \textsuperscript{(89)} ماسكة بإيدها بنت يتيمة
كانت راجعة لديرها.

%%K t14-tit-10 sub
\VUpk{10.2pt}{40}{في طريق الرجوع للبيت.}
\VUsaut{3.0mm}

%%K t14-20-1 p
20. --- وفي اللحظة اللي عمّي رجع فيها لنا، يوحنا ضحك بصوت عالي، وهو شايف الوش
الوسخ واللبس الغريب بتاع اتنين من \VUgras{بتوع المداخن} \textsuperscript{(91)}
المغطّيين بالسناج.

%%K t14-21-1 p
21. --- وكنت عايز أشتري من \VUgras{بايعة الورد} \textsuperscript{(92)} نبتة لأمي، بس
عمّي نبّهني إنها هتبقى تقيلة أوي في الشيل وإن الأحسن نجيب \VUgras{برتقان}. وقرّبنا
من \VUgras{البايعة} \textsuperscript{(94)}، ولما \VUgras{المربّية} \textsuperscript{(93)}
اللي كانت بتختار لتلميذتها خلصت شراها، خلّيناها تخدمنا.

%%K t14-22-1 p
22. --- ولما كنّا راجعين البيت، قابلنا كام واحد نعرفهم: صاحبة قديمة للعيلة، اللي
كانت متأبّطة ذراع \VUgras{وصيفتها} \textsuperscript{(95)}؛ و\VUgras{المهندس}
\textsuperscript{(96)} بتاع الكباري والطرق؛ وواحدة من بنات عمّي مع \VUgras{الدادة}
\textsuperscript{(98)} بتاعتها.

%%K t14-22-2 p
وبعدين قابلنا \VUgras{بتاعة البرانيط} \textsuperscript{(97)} بتاعة أمي، واتنين
\VUgras{رهبان} \textsuperscript{(99)} غربا عن المدينة، اللي جم يسألوا عمّي عن طريق
المحطة.

\VUsaut{6.0mm}
\VUfilet{22mm}
